\verseline{}{你也一样，卡伊埃塔，埃涅阿斯的乳母，临终时}
\verseline{}{把永恒的名声留给了我们的海岸；\footnote{卡伊埃塔是埃涅阿斯的乳母。第六卷结尾船队已停在卡伊埃塔港，本卷随即说明乳母死后葬于此处，海岸因此得名。诗人以第二人称呼告死者； “如果这也算荣耀”给地名不朽蒙上一层克制的迟疑。}}
\verseline{}{如今你的荣誉仍守护此地，你的名字}
\verseline{}{也在辽阔的赫斯珀里亚标记着遗骨——如果这也算荣耀。}
\verseline{5}{而虔敬的埃涅阿斯依礼完成葬仪，}
\verseline{}{筑好坟冢；待深海平息，}
\verseline{}{便张帆继续航程，离开港湾。}
\verseline{}{微风吹入夜色，皎洁的月亮不拒绝}
\verseline{}{照临他们的行路；海面在颤动的光下闪耀。}
\verseline{10}{船只贴着喀耳刻土地近旁的海岸驶过；\footnote{喀耳刻是太阳神之女，以药草把人变作野兽。特洛伊人夜航时只听见受缚狮子、野猪、熊与狼的声音，随涅普顿送来的顺风绕过港湾。战争篇开始前的最后一段海上夜色，仍笼罩着被魔法变形的人声。}}
\verseline{}{在那里，太阳神富有的女儿让人迹不至的林苑}
\verseline{}{回响着不绝的歌声，又在高大的宫室中}
\verseline{}{焚烧芬芳的雪松，照亮夜晚，}
\verseline{}{以鸣响的梭杼穿行在纤细的织线间。}
\verseline{15}{那里传来狮子的低吼和怒声；}
\verseline{}{它们抗拒锁链，直到深夜仍在咆哮；}
\verseline{}{鬃毛森然的野猪、圈栏里的熊}
\verseline{}{也在狂躁，巨狼般的形体也发出嗥叫——}
\verseline{}{凶厉的女神喀耳刻曾以强力药草，把那些人的面容}
\verseline{20}{变作野兽的脸和躯干。}
\verseline{}{为使尽责的特洛伊人不遭这般异变，}
\verseline{}{不被带进那港湾、踏上凶险的海岸，}
\verseline{}{涅普顿以顺风鼓满船帆，}
\verseline{}{给他们逃离之路，载他们驶过沸腾的浅滩。}
\verseline{25}{海面已被霞光染红；在高高的天空，}
\verseline{}{金黄的奥罗拉乘玫瑰色双马车放出光辉；}
\verseline{}{这时风都止息，一切气流忽然}
\verseline{}{沉落，船桨在凝滞如大理石的海上奋力划动。}
\verseline{}{就在此刻，埃涅阿斯从海上望见}
\verseline{30}{一片巨大的树林；秀美的台伯河从林间穿出，\footnote{台伯河在本卷既是航行终点，也是未来定居的地理凭证。原诗写河水裹着金黄泥沙冲入海中，周围有百鸟、歌声和林荫。}}
\verseline{}{急涡翻卷，裹着大量金黄的泥沙，}
\verseline{}{直冲入海。河流四周与上空，}
\verseline{}{惯常栖息岸边和河床的百鸟}
\verseline{}{以歌声抚慰长空，在林中飞翔。}
\verseline{35}{他命同伴转变航向，把船首朝向陆地，}
\verseline{}{欣然驶入林荫笼罩的河流。}
\verseline{}{现在来吧，厄拉托：当时有哪些国王，时世如何，\footnote{“现在来吧，厄拉托”直接呼应阿波罗尼俄斯《阿尔戈英雄纪》第三卷开头对厄拉托的祈请；厄拉托主管爱情诗。这里又是全诗的“第二开篇”：前六卷的漂泊、回忆与入冥结束，本卷起转入意大利战争。}}
\verseline{}{古老的拉提乌姆又是什么局面——当外来的军队}
\verseline{}{带着船队首次抵达奥索尼亚海岸——}
\verseline{40}{我要一一陈说，重新追溯最初战斗的开端。}
\verseline{}{女神，请你启示诗人。我要讲述凶险的战争，}
\verseline{}{讲述阵列，讲述心意受激、奔赴死亡的诸王，}
\verseline{}{讲述伊特鲁里亚军旅，以及整个赫斯珀里亚}
\verseline{}{如何被迫拿起武器。更浩大的世事次第向我展开，}
\verseline{45}{我着手更大的事业。年老的拉提努斯王\footnote{拉提努斯的祖先依次是法乌努斯、皮库斯与萨图尔努斯，因此他的王族血统直接上溯拉提乌姆本地神祇。其唯一儿子早夭，拉维尼娅因而不仅是待嫁女儿，也是继承与王国联盟的中心。}}
\verseline{}{长久在和平中治理安宁的田野与城邑。}
\verseline{}{传说他生于法乌努斯和劳伦图姆宁芙玛丽卡；}
\verseline{}{法乌努斯的父亲是皮库斯，而皮库斯又称你，}
\verseline{}{萨图尔努斯，为他的父亲：你是这血统最初的源头。}
\verseline{50}{依众神的命数，拉提努斯没有儿子，}
\verseline{}{唯一的男嗣刚到青春初年便被夺走。}
\verseline{}{只剩一个女儿维系这家族和广大的宫室；}
\verseline{}{她已经成人，年岁已足，可以婚嫁。}
\verseline{}{辽阔的拉提乌姆和整个奥索尼亚，许多人}
\verseline{55}{都来求娶她；其中最出众、最俊美的}
\verseline{}{是图尔努斯，凭祖先与远祖而有权势；王后\footnote{图尔努斯是鲁图利人的首领，也是拉维尼娅最显赫的本土求婚者；王后阿玛塔支持这桩婚事，视他为理想的女婿。}}
\verseline{}{怀着异乎寻常的偏爱，急于招他为婿。}
\verseline{}{但众神以种种令人惊惧的征兆阻挠此事。}
\verseline{}{宫殿中央幽深的内院里生着一棵月桂树，}
\verseline{60}{枝叶神圣，多年来始终受敬畏守护；}
\verseline{}{相传拉提努斯王初建城堡时发现了它，}
\verseline{}{亲自把它奉献给福玻斯，}
\verseline{}{又由它给劳伦图姆的居民定下族名。}
\verseline{}{密集的蜂群——说来惊异——飞越清澈的天空，}
\verseline{65}{带着巨大的嗡鸣，}
\verseline{}{占据月桂树梢；它们彼此勾连着足，}
\verseline{}{骤然聚成蜂团，从繁叶的枝条垂下。}
\verseline{}{先知立刻说道：“我们看见一位外乡男子}
\verseline{}{正向这里来；一支军队也从同一方向}
\verseline{70}{驶向同一所在，并将统治这最高的城堡。”}
\verseline{}{又有一次，少女拉维尼娅站在父亲身旁，\footnote{拉维尼娅祭坛旁的长发、冠冕与佩饰似乎燃烧，她本人却没有被烧伤。征兆兼具双重意义：少女将因命运与声名显赫，而她的民族将遭逢大战；婚炬、王权光辉与公共灾祸由同一团火连接。}}
\verseline{}{以清洁的火炬燃起祭坛，}
\verseline{}{人们看见火焰攫住她的长发——骇人的异象——}
\verseline{}{她全部佩饰似乎也在噼啪的烈焰中燃烧；}
\verseline{75}{王者的发绺着了火，嵌满宝石的冠冕也着了火；}
\verseline{}{她被烟雾与黄赤的光芒笼罩，}
\verseline{}{把武尔坎努斯的火星撒遍整座宫殿。}
\verseline{}{这景象既可怖，又令人惊异；}
\verseline{}{先知们唱言，她本人将以声名和命运显赫，}
\verseline{80}{却给人民预示一场大战。}
\verseline{}{国王为这些异象忧虑，便去求问能预知未来的}
\verseline{}{父神法乌努斯；他走入高大的林苑，}
\verseline{}{求示于阿尔布涅亚——那是诸林中最广大的一处，\footnote{阿尔布涅亚是法乌努斯神谕所在的圣林与泉。求问者献祭绵羊，夜间卧在羊皮上入睡，借幻象与声音求取答复；这种仪式常被称为孵梦。泉水呼出的有害气息又让这里与地下阿刻戎相连。}}
\verseline{}{圣泉在那里鸣响，幽暗处呼出恶浊的毒气。}
\verseline{85}{意大利各族和整个俄诺特里亚大地}
\verseline{}{遇到疑难都来这里求取答复；祭司把供物带到这里，}
\verseline{}{便在寂静的夜里铺开祭牲羊皮，}
\verseline{}{卧在上面，求取睡梦，}
\verseline{}{便看见众多幻影以奇异的形态飞舞，}
\verseline{90}{听见种种声音，与诸神交谈，}
\verseline{}{又从阿韦尔努斯深处向阿刻戎发问。}
\verseline{}{这一次，拉提努斯父王亲自前来求取答复，}
\verseline{}{依礼宰杀一百只有毛的两齿绵羊，}
\verseline{}{又以它们的皮垫起身体，卧在铺展的羊毛上；}
\verseline{95}{忽然，高处的林中传来声音：}
\verseline{}{“我的后裔，不要寻求把女儿嫁给拉丁男子，}
\verseline{}{也不要信赖已经备好的婚床。}
\verseline{}{外来的女婿将到来，他们的血脉会把我们的名号}
\verseline{}{举向群星；他们一族的后人将看见，}
\verseline{100}{太阳往返、照望两方大洋之间的一切，}
\verseline{}{都转动在自己脚下，受自己统治。”\footnote{法乌努斯禁止拉提努斯把女儿嫁给本族男子，预告外来女婿的后裔将把拉丁名号举向群星。神谕使埃涅阿斯与拉维尼娅的婚姻同时具有家族、王位和未来帝国的意义，也使拉提努斯后来主动提出和平具有充分依据。}}
\verseline{}{拉提努斯没有把父亲法乌努斯的答复}
\verseline{}{以及寂静夜里得到的告诫锁在自己口中；}
\verseline{}{飞旋四方的传闻早已把它带遍}
\verseline{105}{奥索尼亚诸城。这时，拉俄墨冬的特洛伊青年}
\verseline{}{正把船队系在长草的河岸土堤上。}
\verseline{}{埃涅阿斯、主要将领和俊美的尤卢斯}
\verseline{}{在一棵高树的枝荫下歇身，}
\verseline{}{摆开宴食，把小麦烤成的面饼铺在草上，}
\verseline{110}{承托食物——正如朱庇特曾经指示的那样——}
\verseline{}{又把野果堆在这谷物做成的底面上。}
\verseline{}{别的食物恰好吃尽，匮乏迫使他们}
\verseline{}{把牙齿转向那一点点面饼，}
\verseline{}{以手和大胆的牙齿破开命定的圆饼，}
\verseline{115}{连宽阔的饼块也不放过。}
\verseline{}{“嘿，我们连桌子都吃了！”尤卢斯说，\footnote{面饼最初只是承托野果与食物的底面；食物吃尽后，特洛伊人把饼也吃掉，尤卢斯便玩笑说“连桌子都吃了”。少年一句玩笑使先前可怕的预言在日常宴食中应验。}}
\verseline{}{只当一句玩笑，没有再说。这句入耳的话}
\verseline{}{第一次为漫长的苦难带来终结；}
\verseline{}{父亲立刻接住孩子刚出口的话，被神意震住，便止住了他的话头。\footnote{塞尔维乌斯（7.119）将pressit解释为止住阿斯卡尼乌斯的话，好让埃涅阿斯接着祈神；另一解是将孩子的话当作征兆，牢记在心。孩子无意中说出的“吃桌子”，使父亲认出预言已经应验。}}
\verseline{120}{随即说道：“向你致意，命运应许给我的土地！}
\verseline{}{也向你们致意，特洛伊可信的家神！}
\verseline{}{这里是家，这里是祖国。因为父亲安喀塞斯}
\verseline{}{曾把命运的隐秘留给我——如今我才想起：\footnote{第三卷中“吃掉桌子”的预言由刻莱诺说出；到这里，埃涅阿斯却把它归给安喀塞斯。诗中没有解释这一差异。}}
\verseline{}{‘孩子，当你来到陌生的海岸，}
\verseline{125}{食物吃尽，饥饿迫你吃掉桌子，}
\verseline{}{你虽疲惫，也要在那里盼望家园，记住}
\verseline{}{亲手安放最初的屋舍，以土垒筑起居处。’}
\verseline{}{这就是那场饥饿；这就是最后还等着我们的，}
\verseline{}{它将为种种灾祸画下终点。}
\verseline{130}{所以来吧：明天太阳升起第一道光时，}
\verseline{}{让我们欣然查明这是哪里、谁居住于此，}
\verseline{}{他们的城墙又在何处；从港口分道前往探查。}
\verseline{}{现在向朱庇特倾下祭酒，以祷词呼唤}
\verseline{}{父亲安喀塞斯，再把美酒摆上席面。”}
\verseline{135}{说罢，他把带叶的树枝绕在鬓边，}
\verseline{}{向此地的守护神、诸神之首大地母神、}
\verseline{}{宁芙和至今尚不知名的河流祈祷；\footnote{埃涅阿斯依次向地方守护神、大地、宁芙、未知河流、夜与星辰、伊达山的朱庇特、佛里吉亚母神以及天上与地下的双亲祈告；他们此时连当地河流的名字都还不知道。}}
\verseline{}{随后呼唤黑夜与夜里升起的群星，}
\verseline{}{依次呼唤伊达山的朱庇特和佛里吉亚母神，}
\verseline{140}{以及自己的双亲——一位在天，一位在厄瑞玻斯。}
\verseline{}{这时，全能的父神从高天接连三次}
\verseline{}{响亮地发出雷声，又亲手摇动一片}
\verseline{}{燃着光线、闪着金色的云，从天上向他们显现。}
\verseline{}{消息立刻传遍特洛伊队伍：}
\verseline{145}{建立命定城墙的日子已经来临。}
\verseline{}{他们竞相重开宴席，因这重大的吉兆而欢欣，}
\verseline{}{摆下调酒缸，以花冠装点美酒。}
\verseline{}{第二天升起，以第一盏灯火巡照大地；}
\verseline{}{他们分头探查这个民族的城市、疆界和海岸：}
\verseline{150}{这里是努米库斯泉汇成的湖泽，\footnote{努米库斯是拉提乌姆的河流或泉流，后续传统又把埃涅阿斯之死与神化联系到此地。此处它与台伯河、拉丁人共同成为定居地已被辨认的地理证据。}}
\verseline{}{这条是台伯河，强健的拉丁人居住在这里。}
\verseline{}{随后，安喀塞斯之子从各个阶层依次挑选}
\verseline{}{一百名使者，命他们前往国王庄严的城墙，}
\verseline{}{人人都以帕拉斯的枝条覆首，\footnote{百名使者人人以橄榄枝覆首，这是求和者与恳求者的标志。原文 ramis Palladis 以“帕拉斯的枝条”称橄榄枝；此处帕拉斯是弥涅耳瓦的称号，与第八卷以后出现的青年帕拉斯有别。}}
\verseline{155}{向国王献礼，为特乌克罗斯人求取和平。}
\verseline{}{他们毫不迟延，领命后疾步}
\verseline{}{赶路。埃涅阿斯亲自以浅沟划出城墙，}
\verseline{}{整治场地，又把海岸上最初的住所}
\verseline{}{照军营的形制，以城垛和土垒围起。}
\verseline{160}{青年使者已经走完路程，望见拉丁人}
\verseline{}{高耸的塔楼与屋顶，渐渐来到城墙下。}
\verseline{}{城前，少年与初绽青春的年轻人}
\verseline{}{练习骑马，在尘土中驾驭战车，}
\verseline{}{或拉开锐利的强弓，或以臂膀}
\verseline{165}{挥出柔韧的标枪，在竞跑与搏击中相互较量。}
\verseline{}{一名信使策马飞驰，来到年老国王面前，}
\verseline{}{报告身穿陌生服饰的高大男子}
\verseline{}{已经来到。国王命人把他们召进宫内，}
\verseline{}{自己坐在殿堂中央祖传的王座上。}
\verseline{170}{一座高大的庄严宫殿由一百根柱子托起，\footnote{拉提努斯的宫殿同时用于议事、祭祀和宴饮。}}
\verseline{}{坐落在城中最高处，是劳伦图姆皮库斯的王宫；}
\verseline{}{树林和先祖的神圣威严使人敬畏。}
\verseline{}{国王们在这里接过权杖，第一次举起束棒，\footnote{束棒、占卜杖、紫纹礼袍和圣盾都是罗马读者熟悉的权力或祭司标志；维吉尔把它们放进拉提努斯的古老宫廷。}}
\verseline{}{以为吉兆；这里既是他们的议事厅，也是神庙，}
\verseline{175}{又是举行神圣宴礼的席所；宰下公羊后，}
\verseline{}{元老们惯常在绵延的长席旁一同坐下。}
\verseline{}{历代祖先的肖像也依次陈列，}
\verseline{}{以古老雪松雕成：有伊塔卢斯，有父亲萨宾努斯——}
\verseline{}{那位栽培葡萄者的雕像仍握着弯镰——}
\verseline{180}{还有年老的萨图尔努斯和双面雅努斯的形象，}
\verseline{}{都立在门厅；另有从开国以来的诸王，}
\verseline{}{以及为祖国奋战、负过玛尔斯伤痕的英雄。}
\verseline{}{圣洁的门柱上还悬挂许多兵器：}
\verseline{}{被俘的战车、弯曲的战斧、}
\verseline{185}{头盔的冠饰、城门巨大的门闩，}
\verseline{}{标枪、盾牌，还有从船上夺来的撞角。}
\verseline{}{驯马者皮库斯本人也坐在那里，手持奎里努斯的曲柄占卜杖，}
\verseline{}{身束短式紫纹礼袍，左手托着圣盾；}
\verseline{}{他的妻子喀耳刻为情欲所俘，曾用金杖}
\verseline{190}{击打他，以药力改变他的形体，}
\verseline{}{使他变成鸟，又把斑斓颜色洒在双翼上。\footnote{皮库斯是拉提努斯的祖先，也是驯马者与占卜者。原诗称喀耳刻为coniunx（妻子）；塞尔维乌斯（7.190）解作意欲与他成婚者，并述她求爱被拒后把皮库斯变成啄木鸟的传统。}}
\verseline{}{在这般神圣的宫殿和祖先居所中，拉提努斯}
\verseline{}{端坐王位，把特乌克罗斯人召到宫内；}
\verseline{}{待他们入内，他先和颜开口：}
\verseline{195}{“达耳达诺斯人，请说吧——你们的城与族裔}
\verseline{}{我们早已知晓，也听说你们正渡海驶来——}
\verseline{}{你们所求为何？是何缘故，或缺少什么，}
\verseline{}{使船队越过重重碧浪，来到奥索尼亚海岸？}
\verseline{}{无论你们迷失航道，还是受风暴驱迫——}
\verseline{200}{水手在深海常常经受这类苦难——}
\verseline{}{既已进入我们的河岸，在港口停泊，}
\verseline{}{请接受宾主之谊，也请了解拉丁人：}
\verseline{}{我们是萨图尔努斯的民族，凭本心守正，}
\verseline{}{遵循古老神明的习俗，无须枷锁和法律约束。}
\verseline{205}{我还记得——尽管传闻已被岁月遮暗——}
\verseline{}{奥伦库斯的老人们曾说，达耳达诺斯生在这片田野，\footnote{拉提努斯采用达耳达诺斯起源于意大利科律图斯的传统：他从伊特鲁里亚故土迁往佛里吉亚伊达山和萨摩色雷斯，死后列入神位。这样，特洛伊人的“外来”同时又被解释为返回祖先土地。}}
\verseline{}{后来深入佛里吉亚伊达山的城邑，}
\verseline{}{又到色雷斯的萨摩斯，即如今所谓萨摩色雷斯。}
\verseline{}{他正是从这里、从科律图斯的伊特鲁里亚故土出发；}
\verseline{210}{如今，繁星天宫的金色居所把他}
\verseline{}{接上王座，祭坛所敬的神祇又多了一位。”}
\verseline{}{他说完，伊利俄纽斯便接续他的话：}
\verseline{}{“国王，法乌努斯的卓异之子，并非波浪驱赶我们，}
\verseline{}{也非黑沉的风暴迫使我们进入你的土地；}
\verseline{215}{星辰没有使我们偏离方向，海岸也不曾欺骗我们。}
\verseline{}{我们众人经过商议，心甘情愿地}
\verseline{}{来到这座城市；我们被逐出昔日的王国——}
\verseline{}{它曾是太阳从天边的奥林匹斯升起时所见最强盛的一国。}
\verseline{}{我们的族系始于朱庇特；达耳达诺斯一族的青年}
\verseline{220}{以朱庇特为祖先而自豪；我们的国王本人，}
\verseline{}{出自朱庇特最高血脉的特洛伊人埃涅阿斯，派我们来到你的门前。}
\verseline{}{从凶暴的迈锡尼涌向伊达平原的风暴何其广大，}
\verseline{}{又是怎样的命运驱使欧罗巴和亚细亚}
\verseline{}{两方世界彼此冲撞——}
\verseline{225}{连那些住在极远大地、被回流的大洋隔在外面的人，}
\verseline{}{连那些被横亘在四带中央的}
\verseline{}{酷烈日带隔开的人，也全都听说过。\footnote{伊利俄纽斯以欧罗巴与亚细亚的冲突概括特洛伊战争，声称连被大洋隔绝、或隔在中央日带另一侧的人，也已听说此战。古代五带说将世界分作两寒带、两温带，以及中央酷热的日带。}}
\verseline{}{从那场洪流脱身，渡过如此广阔的海面，}
\verseline{}{我们只为祖传诸神求一小块居地、一段}
\verseline{230}{无害的海岸，以及向所有人敞开的水和空气。}
\verseline{}{我们不会使你的王国失色；你们的名声不会因此}
\verseline{}{减轻，这般恩德也不会磨灭；}
\verseline{}{奥索尼亚人绝不会悔于把特洛伊接入怀中。}
\verseline{}{我以埃涅阿斯的命运和有力的右手起誓——}
\verseline{235}{无论谁曾检验过他的信义，还是在兵戈中见识过他：}
\verseline{}{许多民族、许多邦国都曾求取我们——不要因为我们主动}
\verseline{}{手捧圣带、口出祈求之言，就轻视我们——}
\verseline{}{都想把我们接作自己的盟友；}
\verseline{}{但众神以严令驱使我们寻访你们的土地。}
\verseline{240}{达耳达诺斯生于此地；阿波罗又召我们回返，}
\verseline{}{以宏大的神谕催迫我们来到这里，}
\verseline{}{到伊特鲁里亚的台伯河与努米库斯泉的神圣浅滩。}
\verseline{}{此外，埃涅阿斯还献给你昔日家业中的薄礼，}
\verseline{}{那是从燃烧的特洛伊救下的遗物。}
\verseline{245}{父亲安喀塞斯曾用这件黄金器皿在祭坛献酒；}
\verseline{}{普里阿摩斯召集人民、依祖制颁布法令时，}
\verseline{}{曾佩戴这些：权杖、神圣的头冠，}
\verseline{}{以及伊利乌姆妇女织成的衣袍。”\footnote{使者献上的金器、权杖、头冠与衣袍都来自燃烧的特洛伊。它们既是外交礼物，也是王权与祭仪从毁灭中保存下来的遗物；“伊利乌姆妇女的劳作”特指她们织成的衣服。}}
\verseline{}{听完伊利俄纽斯的话，拉提努斯凝住面容，}
\verseline{250}{目光沉向地面，身体不动，只有双眼凝神流转。}
\verseline{}{彩绣的紫袍没有那么打动国王，}
\verseline{}{普里阿摩斯的权杖也一样；他的思绪更多停留在}
\verseline{}{女儿的婚姻和婚床上，}
\verseline{}{在心中反复推求老法乌努斯的神谕：}
\verseline{255}{这就是命运预示的那个人，从异邦而来，}
\verseline{}{受召作女婿，以同等权位共掌王国；}
\verseline{}{他的后代将以卓异的勇力}
\verseline{}{占据整个世界。}
\verseline{}{终于，他欣然说道：“愿众神助成我们的开端，}
\verseline{260}{也助成自己的预兆。特洛伊人，你所求的都会得到，}
\verseline{}{礼物我也不轻慢。只要拉提努斯为王，}
\verseline{}{你们便不缺沃野的丰饶，也不缺特洛伊曾有的富足。}
\verseline{}{只愿埃涅阿斯亲自来——若他如此向往我们，}
\verseline{}{急于缔结宾主之谊、被称为盟友，}
\verseline{265}{就让他前来，不必畏惧友善的面容；}
\verseline{}{对我而言，能与这位王者握手，本身就是和平的一部分。}
\verseline{}{你们则把我的话带回给国王。}
\verseline{}{我有一个女儿；父亲神殿中的神谕不容我}
\verseline{}{把她嫁给本族男子，天空中众多异象也不容许。}
\verseline{270}{预言说，女婿将从外邦海岸而来，}
\verseline{}{这才是留给拉提乌姆的命数；他们将以血脉}
\verseline{}{把我们的名号举向群星。我既相信命运所求正是此人，}
\verseline{}{也愿他正是此人——若我的心还能预见真实。”}
\verseline{}{父王说罢，便从全部马群中挑选骏马；}
\verseline{275}{高高的槽栏中立着三百匹毛色光亮的马。}
\verseline{}{他立刻命人依次给每位特乌克罗斯人}
\verseline{}{牵来飞足的坐骑，铺着紫红和彩绣的毯子；}
\verseline{}{金色胸饰从它们胸前垂下，}
\verseline{}{满身披金，齿间咀嚼着金黄的马衔。}
\verseline{280}{又给不在场的埃涅阿斯一辆战车和一对轭马，}
\verseline{}{它们出自天上的血种，鼻孔喷火；}
\verseline{}{巧艺的喀耳刻曾瞒过父亲，把神骏}
\verseline{}{配给一匹凡间母马，偷育出这些杂血的后代。\footnote{拉提努斯给每位使者一匹马，另给不在场的埃涅阿斯一辆双马战车。驾车神马出自喀耳刻瞒着太阳神父亲、让天上种马与凡间母马交配所得的血统；礼物同时承认埃涅阿斯的首领地位。}}
\verseline{}{埃涅阿斯的使者带着拉提努斯这般赠礼和诺言，}
\verseline{285}{高坐马上返回，带回和平。}
\verseline{}{看啊，朱庇特凶厉的妻子正从伊那科斯的阿尔戈斯归来，}
\verseline{}{乘风而行，停驻空中；}
\verseline{}{她远从西西里帕基努斯上空的高天，望见\footnote{帕基努斯是西西里东南海岬。朱诺从阿尔戈斯返回，在其上空便看见特洛伊人舍船、筑屋，准备定居；他们已不再是一支海上漂泊的船队。}}
\verseline{}{欢欣的埃涅阿斯和达耳达诺斯船队。}
\verseline{290}{她看见他们已开始筑屋，已信赖这片土地，}
\verseline{}{已抛下船只；她停住，尖锐的痛苦把她钉在空中。}
\verseline{}{她摇着头，从胸中倾出这些话：}
\verseline{}{“唉，可恨的族裔！佛里吉亚人的命运，}
\verseline{}{竟总与我们的命运相抗！他们不能倒在西格乌姆平原吗？}
\verseline{295}{被俘后竟不能沦为俘虏？燃烧的特洛伊}
\verseline{}{竟不能烧死这些男子？他们从阵列中央、烈火中央}
\verseline{}{找到了道路。想来，我的神力终究}
\verseline{}{疲惫倒卧，或我已餍足于仇恨，安然止怒。}
\verseline{}{不，我甚至敢在他们被逐出祖国后，仍怀敌意地}
\verseline{300}{追入波涛，在整个海上阻挡这些流亡者！}
\verseline{}{我为对付特乌克罗斯人，耗尽了天空和大海的力量。}
\verseline{}{叙耳提斯、斯库拉、浩大的卡律布狄斯又给我什么助益？\footnote{这段申诉同时叠着两层拉丁前代诗歌的记忆。马克罗比乌斯（6.1.60）把第 295—296 行“被俘后竟不能沦为俘虏？燃烧的特洛伊竟不能烧死这些男子？”与恩尼乌斯《编年史》第十一卷关于特洛伊人在城陷、城焚之后仍未被消灭的句子并列；第 302 行关于叙耳提斯、斯库拉和卡律布狄斯的整组词语，又几乎照取卡图卢斯第 64 首中阿里阿德涅遭忒修斯遗弃后的斥责。}}
\verseline{}{他们藏进渴望已久的台伯河床，}
\verseline{}{不再害怕海洋，也不再害怕我。玛尔斯能够毁灭}
\verseline{305}{凶悍的拉庇泰一族；众神之父本人也把古老的}
\verseline{}{卡吕冬交给狄安娜的怒火；\footnote{朱诺援引两场神怒作比较：玛尔斯曾毁灭拉庇泰人，朱庇特也曾把卡吕冬交给狄安娜惩罚。塞尔维乌斯（7.304、306）将前者联系到婚宴漏请玛尔斯，将后者联系到祭祀遗漏狄安娜；狄安娜遣野猪为害。朱诺用这些先例对照自己的失败。}}
\verseline{}{拉庇泰人犯了什么大罪，卡吕冬又有什么罪，竟配得那般下场？}
\verseline{}{而我，朱庇特伟大的妻子，从未留下任何不敢尝试的事，}
\verseline{}{不幸的我，已经使尽一切手段，}
\verseline{310}{竟败给埃涅阿斯。若我自己的神力}
\verseline{}{还不够强大，我决不迟疑求助世间任何力量；}
\verseline{}{若不能使天神屈服，我便撼动阿刻戎。\footnote{阿刻戎在这里代指冥府力量。朱诺承认自己无法迫使天神改变决定后，转而召来地下的复仇女神。}}
\verseline{}{好吧，我不能阻止他进入拉丁人的王国，}
\verseline{}{命运也已不可动摇地让拉维尼娅成为他的妻子；}
\verseline{315}{可我仍能拖延，给如此重大的事情添上阻碍，}
\verseline{}{仍能毁掉两位国王的人民。}
\verseline{}{就让岳父和女婿用亲人的这般代价结合：}
\verseline{}{少女，特洛伊与鲁图利人的鲜血将作你的嫁妆，}
\verseline{}{贝罗娜等着为你主持婚礼。\footnote{朱诺承认不能阻止埃涅阿斯进入拉丁王国、迎娶拉维尼娅，却仍要拖延并让双方流血。贝罗娜是战争女神，朱诺让她主持婚礼，以特洛伊与鲁图利人的血作嫁妆。基修斯之女是赫库芭，普里阿摩斯的王后；传说她怀帕里斯时梦见生下火炬，预示特洛伊焚毁。朱诺把埃涅阿斯比作另一个帕里斯，以致命婚炬诅咒重建的特洛伊。}不只是基修斯之女}
\verseline{320}{曾怀着一支火炬，生下婚姻的烈火；}
\verseline{}{维纳斯的产儿也一样，是另一个帕里斯，}
\verseline{}{重建的帕加马也将再度燃起致命的婚炬。”}
\verseline{}{说完这些，可怖的女神便降向大地；}
\verseline{}{她从凶神的居所、从地下黑暗召来}
\verseline{325}{制造悲痛的阿勒克托；惨烈的战争、\footnote{阿勒克托是复仇女神之一。朱诺命她撕裂和平、颠覆家庭、使同心者交战并制造战争缘由。阿玛塔本已偏爱图尔努斯；阿勒克托激起她的狂怒，又以婚约与继承权的争端煽动图尔努斯。随后她驱使猎犬追逐家鹿，阿斯卡尼乌斯射伤鹿，引出乡人与特洛伊人的交战。}}
\verseline{}{愤怒、伏击与致命的罪行都合她心意。}
\verseline{}{连父神普路同也憎恶这怪物，她在塔耳塔洛斯的姊妹们也憎恶她：\footnote{原诗此处用 Pluton 称冥界之主，第六卷更常用 Dis；二者都可指同一神祇，这里换用了前者。}}
\verseline{}{她能变出无数面孔，}
\verseline{}{张张凶厉；无数黑蛇从她身上滋生。}
\verseline{330}{朱诺以这些话激起她，开口说道：}
\verseline{}{“黑夜之女，替我专程担起这场劳苦，}
\verseline{}{做成这件事：别让我的荣誉失位，也别让名声受挫}
\verseline{}{而退让；不要让埃涅阿斯的族人以婚姻笼络}
\verseline{}{拉提努斯，也不要让他们盘踞意大利疆土。}
\verseline{335}{你能武装原本同心的兄弟，使他们交战，}
\verseline{}{能以仇恨颠覆家庭，把鞭子与丧葬的火炬}
\verseline{}{送进屋宅；你有千种名字，}
\verseline{}{千种害人的技艺。激荡你那充满恶计的胸怀，}
\verseline{}{撕碎已经缔结的和平，播下战争的罪由：}
\verseline{340}{让青年同时渴求、索要并夺取武器。”}
\verseline{}{于是，浸透戈耳工毒液的阿勒克托}
\verseline{}{首先来到拉提乌姆，来到劳伦图姆君王}
\verseline{}{高耸的宫殿，无声占据阿玛塔的门槛。}
\verseline{}{特乌克罗斯人的到来和图尔努斯的婚事}
\verseline{345}{使她被母亲的忧虑和愤怒烧灼煎熬。}
\verseline{}{女神从深蓝发间投下一条蛇，}
\verseline{}{送进她怀里，直抵胸腹最深处，}
\verseline{}{要使她受这怪物驱狂，搅乱整座王宫。}
\verseline{}{蛇滑入衣袍与光洁的胸前，}
\verseline{350}{盘旋着，不留下触感，瞒过正走向狂乱的王后，}
\verseline{}{把蝮蛇般的气息吹进她胸中；巨蛇化作}
\verseline{}{颈间盘曲的金饰，又化成长圣带垂下的飘带，}
\verseline{}{缠入她的头发，湿滑地游过四肢。}
\verseline{}{最初的瘟毒带着湿润毒液潜入，}
\verseline{355}{试探她的感官，把火缠入骨中；}
\verseline{}{她的心还未完全感到胸中的烈焰时，}
\verseline{}{她的话仍较柔和，仍像母亲往常那样，}
\verseline{}{为女儿与佛里吉亚人的婚姻不断流泪：}
\verseline{}{“父亲，拉维尼娅竟要交给流亡的特乌克罗斯人带走吗？}
\verseline{360}{你既不可怜女儿，也不可怜自己？}
\verseline{}{你也不怜惜她的母亲——第一阵北风一来，}
\verseline{}{那背信的强盗便会掳走少女，驶向大海，把母亲抛下吗？}
\verseline{}{佛里吉亚牧人不正是这样深入拉刻代蒙，}
\verseline{}{把勒达的女儿海伦带往特洛伊城吗？}
\verseline{365}{你的神圣信义何在？你对亲族的旧日关怀何在？}
\verseline{}{你一次次伸给血亲图尔努斯的右手又算什么？}
\verseline{}{如果拉丁人必须从外邦寻求女婿，}
\verseline{}{如果这已定下，父亲法乌努斯的命令又催逼着你，}
\verseline{}{那么我认为，凡是不受我们权杖支配、}
\verseline{370}{与我们分立的自由土地，都是外邦；众神说的也正是如此。}
\verseline{}{至于图尔努斯，若追溯他家族最初的源头，}
\verseline{}{伊那科斯和阿克里西俄斯是他的先祖，家族的根源就在迈锡尼腹地。”\footnote{阿玛塔把埃涅阿斯比作带走海伦的佛里吉亚牧人帕里斯。她把不受拉提努斯权杖支配的土地称为外邦，并追溯图尔努斯的阿尔戈斯祖源，主张他也符合“外来女婿”的神谕。}}
\verseline{}{她以这些话试过拉提努斯，徒劳无功，}
\verseline{}{见他仍不肯让步；毒蛇的狂烈邪力已深入内脏，}
\verseline{375}{游遍全身；}
\verseline{}{这时，不幸的王后才在巨大怪力下暴起，}
\verseline{}{神力扰乱了她，她失去常态，狂奔遍全城。}
\verseline{}{仿佛陀螺有时在拧紧的鞭子下飞旋：}
\verseline{}{孩子们专心游戏，围着空阔庭院站成大圈，}
\verseline{380}{不断抽打它；陀螺受鞭绳驱使，}
\verseline{}{沿弯曲轨迹疾转；上方那群不懂事的}
\verseline{}{孩子惊看黄杨木旋动，}
\verseline{}{每一鞭又给它增添劲力。阿玛塔不比它迟缓，}
\verseline{}{被驱赶着穿过城邑中心和凶悍的人群。}
\verseline{385}{她甚至假托巴克科斯的神力，飞入树林，\footnote{本卷第 385—405 行中，阿玛塔假借巴克科斯狂仪，把拉维尼娅藏进山林，并召集拉丁母亲解开圣带、披兽皮、持缠藤长矛。这组场面重写欧里庇得斯《酒神的女信徒》：狄俄尼索斯使底比斯妇女离城上山，进入集体酒神狂态；阿玛塔同样带领妇女离开家园、进入山林，并使用酒神仪式的服饰和器物。}}
\verseline{}{开始更大的渎神行为，掀起更大的狂乱；}
\verseline{}{她把女儿藏进枝叶繁茂的山间，}
\verseline{}{要从特乌克罗斯人手里夺回婚床，拖延婚炬。}
\verseline{}{她呼喊：“欧吼，巴克科斯！”宣称只有你配得上少女；}
\verseline{390}{因为少女正为你执起柔韧的酒神杖，}
\verseline{}{随舞队绕你起舞，为你蓄养神圣的长发。}
\verseline{}{传闻飞散，同一股狂热驱使所有}
\verseline{}{胸中燃起狂态的母亲去寻觅新的居处。}
\verseline{}{她们弃下家园，把颈项和头发交给风；}
\verseline{395}{另一些以颤动的长啸充满天空，}
\verseline{}{身披兽皮，手执缠着葡萄叶的长矛。}
\verseline{}{王后在众人中央狂热地高举}
\verseline{}{燃烧的松枝，唱起女儿与图尔努斯的婚歌；}
\verseline{}{她转动血红的眼睛，忽然厉声叫道：}
\verseline{400}{“来吧，拉丁人的母亲们，无论身在何处，都听我说：}
\verseline{}{若你们忠诚的心中还留着对不幸的阿玛塔的情分，}
\verseline{}{若对母亲权利的牵挂仍啮咬着你们，}
\verseline{}{就解开束发的圣带，同我领受狂仪！”}
\verseline{}{阿勒克托便在林间、在野兽出没的荒野，}
\verseline{405}{以巴克科斯的尖刺从四面驱策王后。}
\verseline{}{当她觉得自己已经充分挑起最初的狂怒，}
\verseline{}{颠覆拉提努斯的决断和整个家族，}
\verseline{}{阴郁的女神便立即从这里展开暗色双翼，}
\verseline{}{飞向勇敢鲁图利人的城墙；传说那座城市}
\verseline{410}{由阿克里西俄斯之女达那厄率阿尔戈斯移民建立，\footnote{阿尔代亚是鲁图利人的中心。传说达那厄被风带到意大利后率阿尔戈斯移民建城；诗人又说这座城市在自己的时代只余盛名。阿勒克托在此化身朱诺神庙的年老祭司卡吕柏，以熟悉身份接近图尔努斯。}}
\verseline{}{她被疾猛南风带到这里。祖先曾称此地为阿尔代亚，}
\verseline{}{如今“阿尔代亚”的大名仍在，}
\verseline{}{往日的繁盛却已不存。图尔努斯正在城中高屋，}
\verseline{}{在黑夜正中享受睡眠。}
\verseline{415}{阿勒克托卸去凶厉的面孔和复仇女神的肢体，}
\verseline{}{把自己变作老妇的模样；}
\verseline{}{以皱纹犁开狰狞的额头，戴上白发，}
\verseline{}{再以圣带束发，缠入一枝橄榄；}
\verseline{}{她成为朱诺神庙的老祭司卡吕柏，}
\verseline{420}{出现在青年眼前，对他说道：}
\verseline{}{“图尔努斯，你甘愿以往的万般劳苦全都白费，}
\verseline{}{自己的权杖转给达耳达诺斯移民吗？}
\verseline{}{国王拒绝给你婚姻和以鲜血赢来的嫁妆，}
\verseline{}{反倒为王国寻找一个外来的继承人。}
\verseline{425}{去吧，受人嘲笑，再把自己投进得不到回报的险境；}
\verseline{}{去摧毁伊特鲁里亚阵列，以和平庇护拉丁人。}
\verseline{}{你静卧于宁静夜色时，全能的萨图尔努斯之女}
\verseline{}{亲自命我把这些话当面告诉你。}
\verseline{}{所以来吧，欣然准备武装青年，让他们出城；}
\verseline{430}{那些佛里吉亚首领已经驻在秀美的河岸，}
\verseline{}{去焚烧他们和那些彩绘的船只。}
\verseline{}{天神强大的意志如此命令。拉提努斯王本人}
\verseline{}{若不肯答应把女儿嫁给你、履行诺言，}
\verseline{}{就让他终于亲身感受、见识图尔努斯的武力。”}
\verseline{435}{青年嘲笑这位先知，反过来开口回答：}
\verseline{}{“船队驶入台伯河波涛的消息}
\verseline{}{早已传到我耳中，我没有像你猜想的那样毫不知情。}
\verseline{}{不要替我捏造这般恐惧；王者的朱诺}
\verseline{}{也没有忘记我。}
\verseline{440}{母亲，是衰败已久、无力辨真的暮年}
\verseline{}{以无谓的忧虑折磨你；在国王们的战事之间，}
\verseline{}{它又用虚假的恐惧戏弄一个女先知。}
\verseline{}{你该关心的是守护神像和神庙；}
\verseline{}{战与和由男子操持；需要发动战争的是他们。”}
\verseline{445}{这番话使阿勒克托怒火陡燃；}
\verseline{}{青年还在说，突来的战栗已攫住他的四肢，}
\verseline{}{双眼僵直：复仇女神的蛇群无数，齐声嘶鸣，}
\verseline{}{巨大的真容显露；她转动火焰般的眼睛，}
\verseline{}{把仍在犹疑、还想多说的图尔努斯逼退，}
\verseline{450}{从发间竖起成双的蛇，}
\verseline{}{鸣响长鞭，以狂烈的口吻补道：}
\verseline{}{“看吧，我就是那个衰败、被暮年耗尽了判断力，}
\verseline{}{在国王的武装纷争间受虚假恐惧戏弄的人。}
\verseline{}{看清这些：我从凶神姊妹的居所来到这里，}
\verseline{455}{手里带着战争和死亡。”}
\verseline{}{说罢，她向青年掷出火炬，}
\verseline{}{把冒着烟、闪着暗光的火把钉进他胸中。}
\verseline{}{巨大的惊惧撕裂他的睡眠，骨骼和四肢}
\verseline{}{都被全身迸出的汗水浸透。}
\verseline{460}{他神志迷乱，吼着要武器，在床榻和宫室中寻觅武器；}
\verseline{}{对钢铁的渴望与罪恶的战争狂态在他体内翻涌，}
\verseline{}{愤怒更添一层；仿佛柴焰轰然作响，}
\verseline{}{在沸腾铜锅弯曲的腹下越添越旺，}
\verseline{}{水被热力激起，锅内的水流狂涌，}
\verseline{465}{冒着烟、高高溢出泡沫，}
\verseline{}{再不能容于锅中，黑色水汽飞向天空。}
\verseline{}{于是和平已受玷污；他向青年首领下令}
\verseline{}{进军拉提努斯王，又命他们准备武器，}
\verseline{}{守护意大利，把敌人逐出疆界；}
\verseline{470}{他一人足以迎战特乌克罗斯人与拉丁人双方。}
\verseline{}{说完这些，又呼唤众神接纳誓愿，}
\verseline{}{鲁图利人便竞相彼此鼓励，拿起武器：}
\verseline{}{有人被他出众的青春与美貌打动，}
\verseline{}{有人因他祖上皆为国王，有人因他的右手功业显赫。}
\verseline{475}{图尔努斯正以无畏的心气充满鲁图利人的胸膛，}
\verseline{}{阿勒克托则以斯提克斯之翼扑向特乌克罗斯人；}
\verseline{}{她另用一计，察看俊美的尤卢斯正在那段海岸}
\verseline{}{何处设伏、追赶野兽。}
\verseline{}{科库托斯的少女忽然把狂暴投进猎犬，}
\verseline{480}{又以熟悉的气味触动它们的鼻孔，}
\verseline{}{使它们狂热追逐一头牡鹿；这是最初的苦难之因，}
\verseline{}{也在乡人心里点燃了战争。}
\verseline{}{那头牡鹿体态出众，鹿角巨大；\footnote{这头牡鹿由提尔鲁斯一家从幼鹿养大，西尔维娅替它编花、梳毛、洗浴；它习惯人手和餐桌，每晚自行归家。阿勒克托驱使猎犬追鹿，阿斯卡尼乌斯张弓射伤它；鹿带血逃回家中哀鸣。}}
\verseline{}{提尔鲁斯的孩子们从母鹿乳下把它抱走，}
\verseline{485}{与父亲提尔鲁斯一起养大；王家畜群归他管理，}
\verseline{}{广阔田地间看守的职责也托付给他。}
\verseline{}{鹿已惯于听命；孩子们的姊妹西尔维娅倾尽关爱，}
\verseline{}{她以柔软花环编饰鹿角，}
\verseline{}{为野兽梳毛，又在清泉里给它洗浴。}
\verseline{490}{它容许人的手抚摸，习惯主人的餐桌，}
\verseline{}{白日游荡林中，入夜再晚也会}
\verseline{}{自行回到熟悉门槛，走进家里。}
\verseline{}{这天它远远游荡，猎者尤卢斯那群发狂的猎犬}
\verseline{}{惊动了它；它恰好顺着河流}
\verseline{495}{漂向下游，在青绿河岸消解暑热。}
\verseline{}{阿斯卡尼乌斯也燃起对卓异声名的渴望，}
\verseline{}{以弯弓瞄准，射出箭矢；}
\verseline{}{神没有让他的右手射偏；芦箭呼啸，}
\verseline{}{穿过牡鹿腹部和肋下。}
\verseline{500}{受伤的四足兽逃回熟悉的屋宅，}
\verseline{}{呻吟着进入栏圈，浑身带血，发出哀诉，}
\verseline{}{像在恳求，以声音充满整座家园。}
\verseline{}{西尔维娅最先拍着双臂，}
\verseline{}{呼喊救援，召集坚忍的乡人。}
\verseline{505}{他们猝然赶到——因为凶厉的祸害潜伏在寂静林中——}
\verseline{}{一个拿着火烧的棍棒，}
\verseline{}{一个拿着多节的沉重树枝；愤怒把各人寻到的东西}
\verseline{}{都变成武器。提尔鲁斯也召集众人；}
\verseline{}{他原正把楔子打紧，将一棵橡木劈成四块，}
\verseline{510}{此刻抓起斧头，喘息声骇人。}
\verseline{}{凶厉的女神从高处窥望，终于抓住伤人的时机；}
\verseline{}{她飞上畜栏高耸的屋顶，站在最高处，}
\verseline{}{以弯角吹出牧人的信号，}
\verseline{}{把声音拉成塔耳塔洛斯般的长鸣；}
\verseline{515}{整片树林立即战栗，深林回响。}
\verseline{}{远处特里维亚的湖听见了，硫磺水泛白的\footnote{阿勒克托在畜栏屋顶吹响牧人号角，声响传到特里维亚湖、含硫而泛白的纳尔河与韦利努斯泉。原本召集牧人的号角被她吹成塔耳塔洛斯般的声音，迅速变成全境的战争召集。}}
\verseline{}{纳尔河听见了，韦利努斯泉也听见了；}
\verseline{}{惊惧的母亲把孩子紧压在胸前。}
\verseline{}{这时，凶险的号角一声示意，敏捷的乡人}
\verseline{520}{便抓起兵器，从四面奔来；}
\verseline{}{桀骜的农人聚集，特洛伊青年也从敞开的营门}
\verseline{}{涌出，救援阿斯卡尼乌斯。}
\verseline{}{他们排成阵列。如今不再是乡野争斗，}
\verseline{}{不再用硬木棍和烧尖的桩子相击；}
\verseline{525}{双方以双刃铁器决胜，抽出的剑远远}
\verseline{}{森立成乌黑的禾浪；青铜被太阳击中，}
\verseline{}{闪耀，把光芒抛向云下。}
\verseline{}{仿佛海浪刚因初风开始泛白，}
\verseline{}{大海渐渐抬起自身，把波涛举得更高，}
\verseline{530}{最后从深处一直涌向天空。}
\verseline{}{阵线最前，一个青年被啸叫的箭射倒；}
\verseline{}{那是提尔鲁斯诸子中年纪最大的阿尔摩，}
\verseline{}{箭伤留在喉下，鲜血封住了}
\verseline{}{喉中发声的通道，也截断了纤细的生命。\footnote{最先明确倒下的是提尔鲁斯的长子阿尔摩，箭伤喉咙，鲜血堵住声音；随后是前来居间求和的老人伽莱苏斯，他既最富有，也最公正。战争的第一批标志性死者是少年与调停者，英雄决斗尚未开始。}}
\verseline{535}{许多男子也倒在他周围，其中还有年老的伽莱苏斯；}
\verseline{}{他走到两军之间，要带来和平；}
\verseline{}{他曾是奥索尼亚田野上最公正，也最富有的人：}
\verseline{}{五群羊为他鸣叫，五群牛归来，}
\verseline{}{一百把犁替他翻动土地。}
\verseline{540}{双方正在田野势均力敌地交战，}
\verseline{}{女神已经兑现承诺，以鲜血浸透战争，}
\verseline{}{把死亡送进第一场战斗；}
\verseline{}{于是她离开赫斯珀里亚，转身向天，穿过高空，\footnote{Conte 2019 正文取 caelo conuersa（转身向天）。据 该版校勘记，caelo 来自 Schaper 的校改，conuersa 有传本依据；该读法写阿勒克托飞向朱诺的天上居所。传本另有 caeli conuexa 等读法。}}
\verseline{}{以胜利者的傲慢声音对朱诺说：}
\verseline{545}{“看，不和已经为你长成一场悲惨的战争；}
\verseline{}{现在叫他们聚作朋友、缔结盟约吧。}
\verseline{}{既然我已使特乌克罗斯人沾上奥索尼亚的血，}
\verseline{}{只要你的意愿确切，我还会再加上一件：}
\verseline{}{我要以传闻把邻近城市带入战争，}
\verseline{550}{以疯狂的战欲点燃人心，}
\verseline{}{使他们从四方前来支援；我要把武器撒遍田野。”}
\verseline{}{朱诺却回答：“恐惧与欺骗已经足够。}
\verseline{}{战争的缘由已经立稳，双方正在近身交兵，}
\verseline{}{新血已经浸染他们偶然抓到的第一批兵器。}
\verseline{555}{就让维纳斯的卓异之子和拉提努斯王本人}
\verseline{}{庆祝这般婚姻、这般婚礼。}
\verseline{}{至高奥林匹斯的统治者——那位父神——}
\verseline{}{不会愿你更肆意地游荡于天上空气。}
\verseline{}{离开这里；若这些劳苦还有什么后续变故，}
\verseline{560}{我会亲自掌理。”萨图尔努斯之女说完这些，}
\verseline{}{阿勒克托便抬起蛇鸣的翅膀，}
\verseline{}{离开高处，飞往科库托斯的居所。}
\verseline{}{意大利中央的高山之下，有一个地方，}
\verseline{}{名声远播，在许多海岸都有人传说；}
\verseline{565}{那是阿姆普桑克图斯谷。浓密林木以暗色枝叶\footnote{阿姆普桑克图斯谷在今意大利坎帕尼亚，湖泽中涌出含硫气体。诗把它写成狄斯的气孔和裂开的阿刻戎入口，阿勒克托由此返回地下。}}
\verseline{}{从两侧紧逼山谷；谷底，激流回旋，}
\verseline{}{在岩石间轰响。}
\verseline{}{那里显出可怖洞穴和凶厉狄斯的气孔，}
\verseline{}{阿刻戎破土裂开，一道巨大深渊}
\verseline{570}{张开散播瘟毒的咽喉；可恨的复仇女神藏入其中，}
\verseline{}{她的神力一走，大地和天空都卸下重负。}
\verseline{}{与此同时，萨图尔努斯之女、众神之后}
\verseline{}{仍为战争完成最后一步。牧人队伍从战场}
\verseline{}{尽数涌进城内，带回死者：}
\verseline{575}{少年阿尔摩和面容污损的伽莱苏斯。}
\verseline{}{他们呼求众神，又恳求拉提努斯。}
\verseline{}{图尔努斯来到，在杀人的控诉和怒火中}
\verseline{}{把恐惧再度加深：特乌克罗斯人被召进王国，}
\verseline{}{佛里吉亚血统将混入拉丁人，而他被赶出门槛。}
\verseline{580}{那些母亲因巴克科斯震荡而在偏僻林间}
\verseline{}{狂舞，她们的亲属也从四面聚来——阿玛塔绝非无名之人——}
\verseline{}{共同纠缠玛尔斯，索求战争。}
\verseline{}{众人立刻索求那不可言说的战争，违背一切征兆，}
\verseline{}{违背众神的命运，被失常的神力推动，}
\verseline{585}{争相围住拉提努斯王宫。}
\verseline{}{他像海中不动的礁石抵抗人群，}
\verseline{}{像海中礁石迎着巨响而来的风暴；}
\verseline{}{无数波浪环绕吠叫，它仍以自身巨体}
\verseline{}{稳稳坚持；四周的礁岩与泡沫覆盖的石块徒然轰鸣，}
\verseline{590}{海藻撞上它的侧面，又被浪水推回。}
\verseline{}{然而他无力压过人群盲目的决断，}
\verseline{}{事态又依凶厉朱诺的意志行进；}
\verseline{}{父王屡屡请众神和空风作证，终于说道：}
\verseline{}{“唉，我们被命运击垮，被风暴带走！}
\verseline{595}{可怜的人啊，你们将亲自以渎神的鲜血偿还这一切。}
\verseline{}{至于你，图尔努斯，大罪与悲惨的刑罚正等着你；}
\verseline{}{你终将向神明奉上祈愿，却已太迟。}
\verseline{}{而我的安息已经备好，归宿就在门前；}
\verseline{}{我失去的只是一场有福的死亡。”\footnote{拉提努斯正确理解神谕，也接纳了特洛伊使者并提出和平，却压不住宫外的人群和朱诺挑起的战意。他预告图尔努斯日后祈神已晚，又说自己临近生命的安息，却失去了一场有福的死亡；随后他退入宫中，不肯亲手开启战争之门，朱诺便代他撞开铁门。}不再多说，}
\verseline{600}{他把自己封闭在宫内，放下政事的缰绳。}
\verseline{}{赫斯珀里亚的拉提乌姆有一项习俗，后来阿尔巴诸城}
\verseline{}{相继奉为神圣；如今，万邦之首罗马仍在遵行：\footnote{战争之门后来与罗马雅努斯神庙制度相联：和平时关闭，开战时由执政官在元老决议后依礼打开。诗先从古拉提乌姆说到阿尔巴和罗马，再回到拉提努斯；拉提努斯拒绝触门，朱诺便亲自撞开铁门，以神力暴力替代公共仪式。}}
\verseline{}{每当他们最初唤醒玛尔斯、准备战斗，}
\verseline{}{无论要亲手把令人泣下的战争带给格塔人，}
\verseline{605}{还是进攻许耳卡尼亚人、阿拉伯人，或向印度推进，}
\verseline{}{追随奥罗拉，又向帕提亚索回军旗，}
\verseline{}{有两扇战争之门——人们这样称呼它们——}
\verseline{}{因祭仪而神圣，因凶暴玛尔斯而令人畏惧；}
\verseline{}{一百根铜闩和永固的铁梁封闭它们，}
\verseline{610}{守门的雅努斯从不离开门槛。}
\verseline{}{一旦元老们定下开战的决议，}
\verseline{}{执政官本人便身披奎里努斯礼袍，依加比式束衣，}
\verseline{}{以显赫装束打开鸣响的门扉；}
\verseline{}{他亲自召唤战斗，其余青年随声响应，}
\verseline{615}{铜号也以低哑的齐鸣表示赞同。}
\verseline{}{当时，人们也要求拉提努斯依此祖制}
\verseline{}{向埃涅阿斯的族人宣战，打开阴惨的门扉。}
\verseline{}{父王不肯触碰，转身逃开}
\verseline{}{这污秽的职责，把自己藏进幽暗。}
\verseline{620}{于是众神之后从天降下，亲手}
\verseline{}{猛推迟滞的门扉；门轴旋转，}
\verseline{}{萨图尔努斯之女撞开铁固的战争门扉。}
\verseline{}{先前沉寂、未受惊动的奥索尼亚，如今燃烧起来；}
\verseline{}{一部分准备徒步走上平原，一部分高坐骏马，}
\verseline{625}{在尘土中狂奔；所有人都索求兵器。}
\verseline{}{有人用厚厚油脂擦亮光滑的盾牌和明亮标枪，}
\verseline{}{有人在磨石上磨利战斧；}
\verseline{}{他们乐于举起军旗，倾听号角声。}
\verseline{}{五座大城已经架起铁砧，重新锻造兵器：\footnote{阿提纳、提布尔、阿尔代亚、克鲁斯图梅里和安特姆奈五城同时锻造兵器。犁铧、镰刀和农具被放下，祖传的剑重新入炉。《农事诗》1.506--508 也写内战中犁失去荣誉、弯镰被熔成硬剑；马克罗比乌斯（6.1.63）又把《埃涅阿斯纪》7.636 与卢克莱修《物性论》5.1293--1294 关于铁剑和青铜镰的段落并列。}}
\verseline{630}{强盛的阿提纳、骄傲的提布尔、}
\verseline{}{阿尔代亚、克鲁斯图梅里和塔楼环立的安特姆奈。}
\verseline{}{他们凿制坚固的头盔，又弯折柳枝，}
\verseline{}{编成盾牌骨架；另一些人把青铜锻成胸甲，}
\verseline{}{或以柔韧白银打制光滑的胫甲。}
\verseline{635}{对犁铧、镰刀与一切农具的珍爱}
\verseline{}{全都退让；他们在炉中重炼祖传的剑。}
\verseline{}{战号已经鸣响，口令传出战争的信号。}
\verseline{}{这个惊急地从屋中抓起头盔，那个把颤抖的马}
\verseline{}{套上轭具；有人拿起盾牌，披上三层金链编成的}
\verseline{640}{胸甲，把可信赖的剑束在腰间。}
\verseline{}{女神们，现在打开赫利孔，唤起歌声：\footnote{诗人再次祈请缪斯，正式展开意大利军队名录；基本形式承自《伊利亚特》2.484 以下的“舰船目录”。维吉尔的名录除兵力与首领外，还写到河流、城邑、山地、家世、服饰、医术与地方传说。}}
\verseline{}{哪些国王受战争激起？哪些阵列跟随各自的首领，}
\verseline{}{填满田野？慈养万物的意大利大地}
\verseline{}{当时以怎样的男子繁盛，又以怎样的兵器燃烧？}
\verseline{645}{女神们，你们记得，也能一一述说；}
\verseline{}{传到我们这里的，只是传闻的一缕微风。}
\verseline{}{首先步入战争的是来自伊特鲁里亚海岸的凶暴者，}
\verseline{}{蔑视众神的墨泽提乌斯；他武装部众。}
\verseline{}{身旁是他的儿子劳苏斯；除劳伦图姆的}
\verseline{650}{图尔努斯以外，没有谁的身形比他更俊美。}
\verseline{}{劳苏斯善驯烈马，也能制服野兽；}
\verseline{}{他率领阿居拉城徒然追随的一千人——}
\verseline{}{他本值得在父亲治下过得更幸福，}
\verseline{}{也值得有一位并非墨泽提乌斯的父亲。\footnote{墨泽提乌斯是流亡的伊特鲁里亚暴君，诗一开始便称他“蔑视众神”。儿子劳苏斯却俊美、勇敢而忠诚；诗人因此说他本该有另一位父亲。第十卷中，劳苏斯会为救父亲而死。}}
\verseline{655}{他们之后，俊美的阿文提努斯乘着棕榈枝标记的战车，}
\verseline{}{在草地上展示获胜的马匹；他生于俊美的赫尔库勒斯，}
\verseline{}{盾上承载父亲的徽记：}
\verseline{}{一百条蛇，以及蛇群缠绕的勒耳那水蛇。}
\verseline{}{女祭司瑞亚曾在阿文提努斯山林中}
\verseline{660}{秘密分娩，把他带到阳光之下；}
\verseline{}{凡人女子与神结合，那是在提林斯英雄得胜、}
\verseline{}{杀死革律翁，抵达劳伦图姆田野以后，}
\verseline{}{他曾在台伯河中洗浴伊比利亚的牛群。\footnote{阿文提努斯是赫尔库勒斯与女祭司瑞亚之子，名字联系罗马的阿文提努斯山。盾上的勒耳那水蛇、狮皮头罩与赫尔库勒斯装束表明父系身份；革律翁的伊比利亚牛群则把他的谱系接到赫尔库勒斯经过意大利的传说。}}
\verseline{}{部下手持重矛和凶险的暗刃走入战争，}
\verseline{665}{以圆刃短剑和萨宾细矛交战。}
\verseline{}{阿文提努斯本人徒步而行，把巨大的狮皮披卷在身上；}
\verseline{}{可怖的鬃毛未经梳理，白牙森然，}
\verseline{}{罩在头上。他就以这身装束走进王宫，}
\verseline{}{形貌粗野，双肩披着赫尔库勒斯的皮氅。}
\verseline{670}{随后，一对兄弟离开提布尔城墙——}
\verseline{}{那座城以他们的兄弟提布尔图斯为名——}
\verseline{}{卡提卢斯和勇猛的科拉斯，阿尔戈斯的青年；}
\verseline{}{他们穿行密集兵器，奔在阵列最前，}
\verseline{}{仿佛两位云生的半人马从高山顶峰}
\verseline{675}{冲下，疾速离开霍摩勒和积雪的俄特律斯；}
\verseline{}{巨大的树林给他们让路，}
\verseline{}{灌木也在轰然巨响中退开。\footnote{卡提卢斯、科拉斯与兄弟提布尔图斯都属于提布尔建城传统。诗把冲在前列的兄弟比作从霍摩勒与俄特律斯雪山疾降的半人马，森林和灌木为他们轰然让路；这一明喻着重表现速度与冲力。}}
\verseline{}{普莱内斯特城的建立者也没有缺席：}
\verseline{}{据历代人相信，武尔坎努斯之子凯库卢斯生在乡野牲群间，}
\verseline{680}{在炉火旁被发现，后来成为国王。}
\verseline{}{广大的乡间军团跟随他：}
\verseline{}{有居住在高耸的普莱内斯特和加比的朱诺田野的男子，}
\verseline{}{有来自寒冷的阿尼俄恩和溪水润泽的}
\verseline{}{赫尔尼基岩地的；富饶的阿纳格尼亚养育一些人，}
\verseline{685}{另一些由阿马塞努斯河神养育。他们并非人人都有兵器，}
\verseline{}{盾牌与战车也不为人人鸣响；大多数抛出}
\verseline{}{暗青色铅丸，另一些手持两支标枪。}
\verseline{}{黄褐狼皮帽覆盖他们的头；}
\verseline{}{左脚赤裸踏地，}
\verseline{690}{另一只脚套着未经鞣制的皮鞋。}
\verseline{}{接着是墨萨波斯，驯马者、涅普顿之子；\footnote{本卷第 691—705 行写墨萨波斯及其军队。墨萨波斯是涅普顿之子，享有不受火与铁器伤害的神异；他召集费斯肯尼亚、法利斯库斯等地久已不习战争的队伍。众人一路歌唱首领，远望如雪白天鹅从水面升成鸟云；诗在这支军队出场时尤其写到他们的歌声。}}
\verseline{}{神律不容任何人以火伤他，也不容铁器将他击倒。}
\verseline{}{他把久已安居、不惯战争的民族与阵列}
\verseline{}{忽然召入军伍，使他们再度操持铁器。}
\verseline{695}{这些阵列来自费斯肯尼亚和埃奎法利斯库斯人；}
\verseline{}{他们占有索拉克忒的堡垒、弗拉维尼亚的田野、}
\verseline{}{基米努斯的湖与山，还有卡佩纳的林苑。}
\verseline{}{他们排成齐整行列，一路歌唱自己的国王，}
\verseline{}{仿佛雪白的天鹅有时穿过清澈云间，}
\verseline{700}{觅食归来，从长颈发出}
\verseline{}{悠长悦耳的曲调；河流与远处的亚细亚沼泽}
\verseline{}{受歌声震动，一同回响。}
\verseline{}{看见这般大军，没有人会以为铜甲阵列}
\verseline{}{正喧腾作响，只会以为一片嘶哑的鸟云\footnote{Conte 2019 校勘记将 misceri 释为 cum strepitu agitari，写列阵士兵奔行时的呼喊与喧响，和天鹅群嘶哑的鸣声相映。}}
\verseline{705}{从深水旋涡被催逼着飞向海岸。}
\verseline{}{看，克劳苏斯出自萨宾人的古老血统，\footnote{克劳苏斯被罗马克劳狄氏族认作萨宾祖先。密集地名把萨宾、台伯河流域与拉提乌姆诸地的军队依次列出；诗人又以冬季海浪和日照下的麦穗比喻人数之多。}}
\verseline{}{率来一支大军；他自己就抵得上一支大军。}
\verseline{}{如今克劳狄部族和氏族由他繁衍，遍布}
\verseline{}{拉提乌姆——那是在罗马把国权的一份授予萨宾人以后。}
\verseline{710}{一同前来的有阿米忒耳努姆的庞大队伍、古老的奎里特斯，}
\verseline{}{厄瑞图姆全部兵众和盛产橄榄的穆图斯卡；}
\verseline{}{有居于诺门图姆城、韦利努斯湖畔罗塞亚田野的，}
\verseline{}{有居于忒特里卡森严岩壁与塞维鲁斯山的，}
\verseline{}{有居于卡斯佩里亚、佛鲁里和希墨拉河畔的；}
\verseline{715}{有饮台伯河与法巴里斯河水的，寒冷的努尔西亚}
\verseline{}{也派来战士；还有奥耳塔军旅和拉丁人，}
\verseline{}{以及不祥名号的阿里亚河从中穿流的人们。}
\verseline{}{他们像利比亚澄亮如大理石的海面上翻滚的波涛那样众多，}
\verseline{}{当凶暴的俄里翁沉入冬季海水；}
\verseline{720}{又像初升日光烘烤下的密集麦穗，}
\verseline{}{或在赫耳穆斯平原，或在吕喀亚金黄的田野。}
\verseline{}{盾牌鸣响，大地被脚步的重击惊动。}
\verseline{}{随后，与阿伽门农有渊源、特洛伊名号的敌人}
\verseline{}{哈莱苏斯把马套上战车，为图尔努斯带来}
\verseline{725}{一千名凶悍的部众：他们有的以耙翻动}
\verseline{}{巴克科斯赐福的马西库斯田野；有的由高山上的}
\verseline{}{奥伦库斯先辈派来；有的邻接西狄基尼人的平原，}
\verseline{}{或离开卡莱斯；有的居住在浅水的沃尔图尔努斯河边；}
\verseline{}{还有同样粗悍的萨提库卢斯人与奥斯坎人的军队。}
\verseline{730}{他们的武器是圆身短矛，惯于把它}
\verseline{}{系在柔韧皮绳上投出；}
\verseline{}{左手由小圆盾遮护，近战则用弯刃剑。}
\verseline{}{你也不会未经唱述便离开我的诗，}
\verseline{}{俄巴卢斯：相传特隆和塞贝图斯河的宁芙生下了你；}
\verseline{735}{当时特隆年事已高，统治着特勒博伊人的卡普里。}
\verseline{}{但儿子不满足于父亲的土地，}
\verseline{}{那时已以权势广泛统治}
\verseline{}{萨拉斯特斯民族和萨尔努斯灌溉的平野，}
\verseline{}{以及占据鲁弗雷、巴图卢姆、刻莱姆奈田野的人，}
\verseline{740}{还有产苹果的阿贝拉城墙俯瞰的民众。}
\verseline{}{他们惯以条顿方式投掷飞矛，}
\verseline{}{头盔是从软木树上剥取的树皮；}
\verseline{}{青铜小盾发亮，青铜剑也在闪耀。}
\verseline{}{多山的涅尔塞派你参加战斗，}
\verseline{745}{乌芬斯；你以名声和顺遂的武运出众。}
\verseline{}{你的民族埃奎库兰人尤为粗厉，}
\verseline{}{住在坚硬的土地上，惯于在密林久猎。}
\verseline{}{他们全副武装地耕作土地，总爱}
\verseline{}{搬运新近的战利品，靠掠夺生活。}
\verseline{750}{一位祭司也从马鲁维亚民族到来，}
\verseline{}{头盔上缀着枝叶和吉祥的橄榄；}
\verseline{}{他奉阿耳喀普斯王之命，是最勇敢的翁布罗。\footnote{翁布罗来自马尔西人的治蛇传统，能以歌声、手势和草药使毒蛇入睡、缓解咬伤，却不能治愈达耳达诺斯枪尖造成的创口。诗随即写安吉提亚林与福基努斯湖为他的死亡哭泣。}}
\verseline{}{面对蛇类和喷吐毒气的水蛇，}
\verseline{}{他惯以歌声和手势把睡眠撒向它们，}
\verseline{755}{抚平它们的怒性，以技艺解除咬伤。}
\verseline{}{然而他无力医治达耳达诺斯枪尖的一击；}
\verseline{}{催眠歌声没能救治他的创伤，}
\verseline{}{从马尔西山中采来的药草也没有。}
\verseline{}{安吉提亚的林苑为你哭泣，福基努斯玻璃般的水波，}
\verseline{760}{清澈的湖泊也为你哭泣。}
\verseline{}{希波吕托斯最俊美的儿子维尔比乌斯也走向战争；\footnote{此处出战的维尔比乌斯是希波吕托斯之子。其父被马拖死后由医神阿斯克勒庇俄斯复活；“重返高天”写从冥界重回生者世界。“派翁尼亚药草”可按医神派翁的药草理解。朱庇特以雷霆击落医神，狄安娜把希波吕托斯藏在阿里喀亚，改名维尔比乌斯。圣林禁马纪念旧日死亡，儿子却仍驾马奔赴战争。}}
\verseline{}{他的母城阿里喀亚派出这位出众的青年，}
\verseline{}{他在湿润岸边的厄革里亚林苑中长大，}
\verseline{}{那里有狄安娜祭礼丰厚、蒙神悦纳的祭坛。}
\verseline{765}{传说希波吕托斯因继母的诡计而死，}
\verseline{}{以自己的鲜血偿尽父亲施加的刑罚；}
\verseline{}{受惊的马把他撕裂，他后来又升向群星，}
\verseline{}{重返高天、天界的气息之中——}
\verseline{}{派翁尼亚的药草与狄安娜的爱把他召回。}
\verseline{770}{全能的父神震怒：竟有一个凡人}
\verseline{}{从地下幽影升回生命之光，}
\verseline{}{他便亲自以雷霆将这般医术的发明者——}
\verseline{}{福玻斯之子——击入斯提克斯水中。\footnote{Conte 2019 第773行取 Phoebigenam（福玻斯所生者），指医神阿斯克勒庇俄斯。该版校勘记记录 Poenigenam 等异读，并认为 Poenigenam 的构词在此不能接受。}}
\verseline{}{慈爱的特里维亚却把希波吕托斯藏进隐秘居所，}
\verseline{775}{把他交给宁芙厄革里亚和她的林苑，}
\verseline{}{使他独自在意大利密林里默默度过一生，}
\verseline{}{又改换名字，成为维尔比乌斯。}
\verseline{}{因此，坚蹄的马也被隔绝在特里维亚的神庙}
\verseline{}{和圣林之外：当年它们受海中怪物惊吓，}
\verseline{780}{在岸边掀翻战车和青年。}
\verseline{}{但他的儿子仍在平展田野上驱驰烈马，}
\verseline{}{乘战车奔向战争。}
\verseline{}{图尔努斯本人身形出众，行在最前，\footnote{图尔努斯高出全军，三重盔鬃上托着喷火的喀迈拉；诗又说，战场流血越多，盔饰上的怪兽就越凶猛。}}
\verseline{}{手持兵器，比众人高出整个头颅。}
\verseline{785}{他的高盔蓄着三重鬃羽，托起一头喀迈拉，}
\verseline{}{咽喉喷出埃特纳山的火焰；}
\verseline{}{战斗越因流血而变得凶厉，}
\verseline{}{怪兽便越发狂吼，放出阴惨的烈火。}
\verseline{}{他光滑的盾面以黄金饰着高举双角的伊俄，\footnote{图尔努斯盾上刻着已经长毛、化为母牛的伊俄，守卫她的百眼阿耳戈斯，以及她的父亲、河神伊那科斯。盾面把图尔努斯所宣称的阿尔戈斯祖源化为醒目的家族图案。}}
\verseline{790}{她已遍身长毛，已经化作母牛——}
\verseline{}{一幅宏大的图案；还有少女的守卫阿耳戈斯，}
\verseline{}{以及父神伊那科斯从雕刻的瓮中倾出河流。}
\verseline{}{步兵如云跟随；持盾的队伍}
\verseline{}{密密铺满整片平原：有阿尔戈斯青年、}
\verseline{795}{奥伦库斯军众、鲁图利人、古老的西坎尼亚人、}
\verseline{}{萨克拉尼阵列和盾牌彩绘的拉比库姆人；}
\verseline{}{有在你的林地、台伯河啊，在努米库斯圣岸}
\verseline{}{耕作的；有以犁翻动鲁图利山丘和}
\verseline{}{喀耳刻山脊的；安克苏鲁斯的朱庇特守护一些人的田野，\footnote{图尔努斯后的地名群从台伯河、努米库斯与鲁图利山丘延伸到喀耳刻山、安克苏鲁斯的朱庇特、费罗尼亚林、萨图拉沼泽和乌芬斯河。这里的军队由被耕作、被神守护并有河流出海的地方组成，具有具体的地理根系。}}
\verseline{800}{费罗尼亚也守护这些田野，在绿荫林苑里欢喜；}
\verseline{}{还有萨图拉黑色沼泽所在之地，寒冷的乌芬斯河}
\verseline{}{从深谷寻找道路、隐没入海之处。}
\verseline{}{最后，卡米拉从沃尔斯基民族到来，\footnote{卡米拉是沃尔斯基人的少女战士，率骑兵与青铜闪耀的队伍压轴入场。她不习纺锤和织篮，却能忍受艰苦战斗；奔跑不伤麦穗、渡海不湿足的夸张速度，把她写得近乎超人。}}
\verseline{}{率领骑兵战阵和闪耀青铜的队伍；}
\verseline{805}{这位女战士的双手不惯于纺锤和弥涅耳瓦的织篮；}
\verseline{}{这少女所习惯的，却是战斗——}
\verseline{}{承受苦战，以双足奔跑胜过疾风。}
\verseline{}{她能飞掠未割禾田的梢头，}
\verseline{}{奔行却不伤柔嫩的麦穗；}
\verseline{810}{也能悬在翻涌的海浪上，穿过海心，}
\verseline{}{疾驰而过，足底不沾海水。}
\verseline{}{所有青年从屋舍与田野涌出，}
\verseline{}{母亲们也聚成一群，惊叹地望着她行进，}
\verseline{}{心神震动，张口凝视：她光洁的双肩}
\verseline{815}{如何披着王者的紫袍；金扣如何以黄金绾发，}
\verseline{}{她如何亲自背着吕喀亚的箭袋，}
\verseline{}{手持顶着枪尖的牧人桃金娘杖。\footnote{青年与母亲们从屋舍、田野出来观看卡米拉，最后的特写依次落在紫袍、金扣、吕喀亚箭袋与带枪尖的牧人桃金娘杖。原诗以武器结束本卷；第八卷开篇立即写图尔努斯举起战争信号，号角鸣响，意大利名录由此转入实际集结。}}
